\documentclass[a4paper,10pt]{article}
\usepackage[utf8]{inputenc}
\RequirePackage{amsthm,amsmath,amssymb,amscd}
\def\D{^{C}D_t^\alpha}
\def\W{\mathbb{W}}
\def\H{\mathbb{H}}
\def\K{\mathbb{K}}
\def\E{\mathbb{E}}
\def\cL{\mathcal{L}}
\def\R{\mathbb{R}}
\def\ud{\text{d}}
\begin{document}

Here are comments for 6d165343 on the file -- {\it thesis.tex}:

\begin{enumerate}
  \item p. 2 the add ``and" after the centered equation for $J_0$.
  \item p. 2, $\Gamma$ is the standard {\it Gamma-function} $\to$ $\Gamma(\cdot)$ is the standard Gamma function.
  \item You may group the last paragraph in a remark.
  \item p.3, line 2: give the two names for the references [MS21]. You use {\it et al} (with italic style) only when
    there are more than three authors.
  \item p. 2, the equation for $\sigma_N(x)$, the word "continuous and" is confusing. Should we just say that
    \begin{align*}
      \sigma_N(x) = \sigma\left(\left( 1 \wedge \frac{N}{|x|}\right )\: x \right).
    \end{align*}
  \item Last paragraph of Section 1, "It\^o" should be properly typed; I would say "the $L^2(\Omega)$ existence" with
    "the".
  % \item Can you extend Lemma 2.2 to include the case when $\beta=2$ (you may think $\alpha>0$)? I should be able to
  % share with you a recent Ph.D. work co-supervised by me in this week or so, where the fundamental solutions are given
  % for all cases $\beta\in (0,2]$, $\alpha>0$ and $\gamma\ge 0$.
  \item Eq. (2.7), I would have $k>0$ and write the second part as
    \begin{align*}
      \sup_{t\ge 0} \int_0^t se^{-a s} \ud s = \int_0^\infty se^{-a s} \ud s = \frac{1}{a^2} \quad \text{for all $a>0$}
    \end{align*}
  \item Regarding Proposition 2.4:
     \begin{enumerate}
       \item Did you forgot subscript $0$ for $J$ in the beginning of the proof? There also misses a period at the end
         of this centered equation. Check all thesis about such typos.
       \item You may put the last two centered equation as one equation on p. 4.
       \item You may remove [MS21, Proposition 3.2] since you are going to prove this one. But you may mention the
         relation of this proposition with Proposition 3.2 of [MS21].
       \item First centered equation in the proof, you forgot the period after $\beta\in(1,2)$
       \item Eq. (2.10), $\sup_{t\ge 0}\sup_{x\in\R}$ should be $\displaystyle\sup_{t\ge 0,\:x\in\R^d}$.
       \item It is cumbersome to put the two cases formula in other equations. Let's define the right-hand side of
         (2.10) by $C_\beta$ and then simplify your presentation.
     \end{enumerate}
   \item Proposition 2.7, put two conditions on $a$ and $p$ in one line (there is an unnecessary period).
   \item Regarding the proof of Theorem 1.1:
     \begin{enumerate}
       \item the Picard iteration should be written with $1$ replaced by $J_0(t,x)$ since you have equality here.
         Moreover, the upper bound of $J_0(t,x)$ is function of $t$ in general even though we can make it uniform in
         $x$. But since we do not consider asymptotics in $t$, for fixed $t\in [0,T]$, we can say that $|J_0(t,x)|\le
         C_T$ so that when we carry out the upper bound moment estimates, there won't be any issues.
       \item The last sentence: This completes the proof of Theorem 1.1.
       \item The rate $\frac{\ln(2N)^{p\delta/2}}{N^p}$ makes one think the result is far from being sharp. we are using
         polynomial to combat against logarithm. This is unfair even this does prove our results. Can you check [MS21]
         and Dalang's original work to see if this is indeed the case? There have to be a lot room to improve such
         result.
       \item The two-line equation for $\mathbb{P}\left(\tau_N<T\right)$ can be made in one line; the next two-line
         equation looks too wide. Try to improve the presentation of the paper.
       \item At beginning of Section 4, always find a proper place to introduce your convention on Fourier transform.
       \item Assumption 4.1, Use period at each end of items and capitalize the initial letter.
     \end{enumerate}
\end{enumerate}

\end{document}
